\clearpage
\section{GNSS-Messungen Schwanderbärgli}
\label{sec:gnss}
Eines der Ziele der Netzmessung 2026 des Deformationsnetzes Schwanderbärgli ist die Modernisierung des bestehenden 
Messkonzeptes. Im aktuell stark auf terrestrisch tachymetrische Messungen ausgelegten Netz soll untersucht werden, 
wie der Messaufwand reduziert werden kann. Konkret sollen lange, gegenseitige Visuren wo möglich weggelassen oder 
durch einseitige Visuren ersetzt werden. Um dennoch eine genügend hohe Redundanz und geometrische Stabilität zu 
erhalten, wird vermehrt auf GNSS-Messungen gesetzt.

\begin{figure}[h!]
    \centering
    \includegraphics[width=\textwidth]{01_Netzübersicht_GNSS-v4.jpg}
    \caption[Bildbeschriftung]{Netzübersicht GNSS-Messung Feldkurs 2026}
    \label{fig:abb-gnss}
\end{figure}


\subsection{Messkampagne 2026}
Abbildung~\ref{fig:abb-gnss} zeigt das überarbeitete Messkonzept für das GNSS Grundlagennetz 2026.
Nebst der Punktauswahl für die langstatischen GNSS-Messungen wurde auch der Messaufbau im Vergleich zu früheren 
Messkampagnen angepasst. Vier der GNSS Punkte (grün in Abbildung~\ref{fig:abb-gnss}) dienen neu als lokale Referenzstationen, 
zu welchen die Basislinienberechnungen der weiteren GNSS-Stationen (gelb) erfolgt. Dies bringt den Vorteil, dass über 
den gesamten vertikalen Projektperimeter Referenzdaten auf ähnlicher Höhe verfügbar sind. 

Die GNSS-Messungen der Messkampagne 2026 erfolgten zwischen Dienstag und Donnerstag, 08.-10. September 2026. Die 
lokalen Referenzstationen (9500, 7015, 7004, 7000) wurden durchgängig während 43~-~47~h gemessen. Währenddessen 
wurden die Verdichtungspunkte in 4~-~22~h Sessionen gemessen. Abbildung~\ref{fig:abb-gnss-zeiten} zeigt 
die zeitliche Verteilung und Überlappungen der verschiedenen GNSS-Stationen. Alle Messungen erfolgten mit insgesamt 
acht Leica GS18i Empfängern der FHNW Muttenz. Die Messungen wurden in einem 15~s Intervall, ohne elevations- oder 
DOP-Abhängigem Cut-off aufgezeichnet.

\begin{figure}[h!]
    \centering
    \includegraphics[width=\textwidth]{GNSS_Zeitfenster.jpg}
    \caption{Zeitliche Einteilung der GNSS-Messungen 2026. Rot die Referenzstationen, blau die Verdichtungspunkte. Abbildung aus \cite{walti_vertiefungsprofil_2026} verwendet.}
    \label{fig:abb-gnss-zeiten}
\end{figure}


\subsection{Datenaufbereitung}
\label{sec:datenaufbereitung}
Die Plausibilitätsprüfung und Kontrolle auf Vollständigkeit erfolgte bereits während dem Feldkurs und ist in 
\cite{walti_vertiefungsprofil_2026} dokumentiert. 

\subsection{Datenauswertung}
\label{sec:datenauswertung}
Die Auswertung der GNSS Daten setzt sich aus zwei Teilen zusammen. In einem ersten Schritt sollen die definitiven 
Koordinaten aller langstatischen GNSS Punkte berechnet werden, welche zur Lagerung der tachymetrischen Messungen 
dienen. In einem zweiten Schritt werden Einflüsse unterschiedlicher Auswertestrategien auf die 
Koordinaten der GNSS-Messungen untersucht.

Die Berechnung aller Basislinien erfolgt in Leica infinity (v5.1.0) \parencite{leica-geosystems_leica_2026}. Die Deformationsmessung 
Schwanderbärgli wird im Bezugssystem CH1903+ (LV95/LHN95) ausgewertet.

\subsubsection{Koordinatenberechnung GNSS Grundlagennetz}
Da für die vier lokalen Referenzstationen nicht ausgeschlossen werden kann, dass sie von Deformationen betroffen sind, 
werden deren Basislinien zur nächstgelegenen AGNES-Station Hasliberg (HABG) berechnet. Diese liegt 
rund 10 km vom Projektperimeter entfernt auf 1089 m.ü.M (CH1903+) \parencite{swisstopo_hfp1_2026}. 

Die Zeitreihendarstellung des GDOP und der 3D-Residuen zur statischen Lösung weisen keine nennenswerten Ausreisser 
auf, weshalb die Messdaten nicht weiter manuell bereinigt werden. Für die Basislinienberechnung werden die präzisen Ephemeriden 
des Center for Orbit Determination in Europe (CODE) verwendet. Die ionosphärischen Korrekturen werden aus den 
Messdaten geschätzt, während zur Modellierung der Troposphäre die Vienna Mapping Function (VMF) 
\parencite{boehm_vienna_2004} mit dem Global Pressure and Temperature Model GPT2 \parencite{lagler_gpt2_2013} 
verwendet werden. Die finalen Koordinaten der Punkte 7000 und 7004 werden noch um den Offset aus der Dreifusszentrierung 
(vgl. hierzu auch \cite{walti_vertiefungsprofil_2026}) korrigiert.

Anschliessend werden die Verdichtungspunkte (gelb in Abbildung \ref{fig:abb-gnss}) zu den lokalen Referenzstationen 
berechnet. Aufgrund erhöhter Abweichungen der Basislinien zur Station 7004 (\todo{verweis auf Vergleichskapitel}) wird 
diese nicht zur finalen Basislinienberechnung verwendet, sondern ebenfalls als Verdichtungspunkt berechnet. 
In der Zeitreihendarstellung der Messungen im Punkt 6011 
sind vermehrt Ausreisser im Meterbereich in den 3D-Residuen sichtbar. Diese sind auf starke Abschattung im Punkt 
und die daraus resultierende schlechte Satellitengeometrie zurückzuführen. Die betreffenden Zeitfenster 
werden manuell bereinigt. Die Positionsresiduen und der GDOP der 
anderen Punkte sind unauffällig. Zur Berechnung werden die gleichen Einstellungen wie zur Berechnung lokalen Referenzstationen 
verwendet.

Pro Punkt ergeben sich somit drei Koordinaten, jeweils zu jeder Basisstation. Diese werden nach der Höhendifferenz 
des Punktes zur entsprechenden Basisstation linear gewichtet gemittelt, wobei die kleinste Höhendifferenz das höchste 
Gewicht erhält. Die gemittelten Koordinaten werden ebenfalls wo nötig um den 
Offset der Dreifusszentrierung korrigiert und anschliessend in ein .mes-File konvertiert. Es werden zwei Files erstellt, 
eines mit ellipsoidischen Höhen zum Vergleich mit der trigonometrischen Höhenbestimmung (THB) und eines mit 
orthometrischen Höhen zur Verwendung als Lagerungspunkte für die tachymetrischen Messungen. Die finalen Koordinaten 
sind in den Anhängen \todo{ANhang referenzieren} beigelegt.

\subsubsection{Einfluss des Koordinatenbezugssystem}
Die Gebrauchshöhen LN02 wurden ausschliesslich aus Nivellementmessungen berechnet. Da weder die Lotabweichung 
noch die lokale Schwere korrigiert wurden, handelt es sich um orthometrische Höhen. Der neue Höhenreferenzrahmen 
LHN95 hingegen basiert sowohl auf Nivellement-, als auch auf GNSS-Messungen. Hierbei wurden die 
Lotabweichung und die lokale Schwere korrigiert, woraus streng orthometrische Höhen resultieren. Während entlang 
der Nivellementlinien die Differenzen zwischen LN02 und LHN95 wenige Millimeter betragen, können diese 
ausserhalb der Linien auf bis zu 65 cm anwachsen \parencite{schlatter_neue_2006}. Zwischen den 
beiden Referenzrahmen besteht nur eine Näherungstransformation (HTRANS), aber keine streng mathematische Beschreibung. 
Die Transformation zwischen den beiden Systemen kann somit zu einem Genauigkeitsverlust führen. Im Schwanderbärgli 
beträgt der Höhenoffset entlang der Fixpunkte im Tal (9000 - 9500 - TP510) rund 2.8 cm. Mit zunehmender Höhe steigt 
die Differenz bis auf > 10 cm Höhendifferenz auf dem Punkt TP506.


\subsubsection{Bereinigen Messungen}
\subsubsection{Basislinienqualität}
\subsubsection{Mittelbildung / Einfluss der Referenzstationen}


\subsubsection{Verschiedene Ionosphären-Modelle}
\label{sec:ionosphäre}
Da die Ionosphäre ein dispersives Medium ist, können die meisten ionosphärischen Effekte durch die unterschiedliche
Verzögerung des Signals (z.B. L1 und L2) eliminiert werden. Dies ist jedoch nur möglich, wenn mit Mehrfrequenzempfängern
gemessen wurde. Infinity rechnet sobald mehr als 45min Messdauer vorhanden sind ein eigenes Modell der Verzögerungen. 
Im Prozess der Auswertung wird untersucht, wie viel Einfluss die unterschiedliche Modellierung der ionosphärischen
Verzögerungen auf die Messresultate hat.

Grundsätzlich herrschte während der Messdauer keine besonderen Verhältnisse in der Ionosphäre, wie gemäss Abbildung~\ref{fig:abb-I95} ein Ausschnitt aus der I95 Dokumentation der Swisstopo zeigt.

\begin{figure}[h!]
    \centering
    \includegraphics[width=\textwidth]{I95-Doku-triple.png}
    \caption{Dokumentation des Ionosphäre gemäss I95 der Swisstopo}
    \label{fig:abb-I95}
\end{figure}

Im ersten Versuch wurde als Ionosphärenmodell eine tägliche Berechnung der Uni Bern (CODE) verwendet (.ION - File).
Diese Lösung zeigt in der Lage deutliche systematische Abweichungen zu unserer ursprünglichen Berechnung,
in welcher der Einfluss mittels Zweifrequenzen eliminiert wurde.
Die Ergebnisse sind in folgender Tabelle~\ref{tab:iono-ber} dokumentiert.
\begin{table}[htbp]
    \centering
    \caption{Differenzen zwischen den Koordinaten. (Ionosphärenmodell - Berechnet)}
    \label{tab:iono-ber}
    \begin{tabular}{rccc}
        \toprule
        \textbf{Pkt.-Nr.} & \textbf{Ost [m]} & \textbf{Nord [m]} & \textbf{Ortho. Höhe [m]} \\
        \midrule
        \rowcolor{lightgray}
        7000 &  0.0067 & -0.0023 &  \textcolor{red}{0.0171} \\
        7004 &  0.0056 & -0.0018 &  0.0095 \\
        7015 &  0.0057 & -0.0019 & -0.0046 \\
        9500 &  0.0059 & -0.0011 &  0.0058 \\
        \bottomrule
    \end{tabular}
\end{table}

Im zweiten Versuch wird das weltweite Ionosphärenmodell von Klobuchar (1996) \parencite{j_sanz_subirana_klobuchar_2011},
verwendet. Dies wurde  hauptsächlich für Einfrequenzempfänger entwickelt und aus empirisch aus Erfahrungswerten bestimmt.
Die Ergebnisse können der Tabelle~\ref{tab:klobuchar-ber} entnommen werden.
\begin{table}[htbp]
    \centering
    \caption{Differenzen zwischen den Koordinaten. (Klobucharmodell - Berechnet)}
    \label{tab:klobuchar-ber}
    \begin{tabular}{rccc}
        \toprule
        \textbf{Pkt.-Nr.} & \textbf{Ost [m]} & \textbf{Nord [m]} & \textbf{Ortho. Höhe [m]} \\
        \midrule
        \rowcolor{lightgray}
        7000 & -0.0040 &  0.0012 &  \textcolor{red}{0.0107} \\
        7004 & -0.0037 &  0.0008 &  0.0010 \\
        7015 & -0.0035 &  0.0004 & -0.0054 \\
        9500 & -0.0033 & -0.0001 &  0.0039 \\
        \bottomrule
    \end{tabular}
\end{table}

Schlussendlich stellt sich heraus, dass die ursprünglich berechnete Version, sich grob in der Mitte der beiden anderen Versionen befindet.
In der Lage sind dabei deutlich grössere systematische Abweichungen erkennbar als in der Höhe. Deshalb kann daraus geschlossen werden,
dass das berechnete Ionosphärenmodell die besten Lösungen erbringt. 


\subsubsection{Verschiedene Troposphären-Modelle}

Die Troposphäre ist kein dispersives Medium, weshalb ihre Verzögerung nicht mittels Mehrfrequenzempfängern eliminiert werden kann.
Wie in Kapitel~\ref{sec:literaturrecherche} behandelt, besteht die Troposphäre aus einem homogenen, gut modellierbaren - sowie einen
hydrostatischen, schlecht modellierbaren Teil. Je nach Elevationswinkel kann die Signalverzögerung bis zu mehreren Metern ausmachen,
was sich stark auf die Positionsgenauigkeit auswirkt. Ziel ist es also, für diese Korrektur das bestmögliche Modell zu verwenden
oder zu berechnen. In Infinity stehen rund fünf verschiedene Methoden als Auswertestrategie für das Troposphärenmodell zur Verfügung.
Die wichtigsten beiden sind «Berechnet» und «VMF mit GPT2-Modell». Die übrigen, namentlich: Essen und Froome (1951),
Hopfield/Modifiziertes Hopfield (1969) und Saastamoinen (1972) sind bereits etwas veraltet und ziehen einen Qualitätsverlust mit sich.
Infinity verwendet für Messungen von weniger als 45min die Vienna Mapping Function und für Messungen darüber berechnet sie ihr
eigenes Troposphärenmodell.

Wie im Kapitel~\ref{sec:ionosphäre} wird die Auswirkung dieser Auswertemethoden auf die Punktgenauigkeit untersucht.
Dabei wird das Resultat des berechneten Modells mit jenem der VMF mit GPT2-Modell verglichen. Tabelle~\ref{tab:vmf-ber} ist zu entnehmen,
dass sich die beiden Punkte in der Lage mit einem Wert von weniger als einem Millimeter kaum unterscheiden. In der Höhe jedoch
variiert der höchste Punkt 7000 wiederum um über einem Zentimeter. Die übrigen Punkte bewegen sich im Rahmen von lediglich 1-5 mm.

\begin{table}[htbp]
    \centering
    \caption{Differenzen zwischen den Koordinaten. (VMF - Berechnet)}
    \label{tab:vmf-ber}
    \begin{tabular}{rccc}
        \toprule
        \textbf{Pkt.-Nr.} & \textbf{Ost [m]} & \textbf{Nord [m]} & \textbf{Ortho. Höhe [m]} \\
        \midrule
        \rowcolor{lightgray}
        7000 &  0.0000 &  0.0000 &  \textcolor{red}{-0.0118} \\
        7004 &  0.0001 &  0.0000 & -0.0016 \\
        7015 &  0.0001 &  0.0001 &  0.0042 \\
        9500 &  0.0001 &  0.0000 & -0.0053 \\
        \bottomrule
    \end{tabular}
\end{table}

Um den Einfluss der alten Modelle besser einschätzen zu können, werden alle Beispielpunkte noch mit dem Saastamoinen Modell berechnet.
Dies wird stellvertretend für alle älteren Modelle verwendet. Gemäss Tabelle~\ref{tab:saas-ber} sind in der Lage hierbei wider minimale systematische
Abweichungen von bis zu einem halben Millimeter zu erkennen. In der Höhe ist das Modell um 2-4 mm auffallend nahe an der VMF.

\begin{table}[htbp]
    \centering
    \caption{Differenzen zwischen den Koordinaten. (Saastamoinen - Berechnet)}
    \label{tab:saas-ber}
    \begin{tabular}{rccc}
        \toprule
        \textbf{Pkt.-Nr.} & \textbf{Ost [m]} & \textbf{Nord [m]} & \textbf{Ortho. Höhe [m]} \\
        \midrule
        \rowcolor{lightgray}
        7000 & -0.0005 &  0.0001 &  \textcolor{red}{-0.0143} \\
        7004 & -0.0004 &  0.0001 &  0.0007 \\
        7015 & -0.0003 &  0.0002 &  0.0075 \\
        9500 & -0.0004 &  0.0001 & -0.0013 \\
        \bottomrule
    \end{tabular}
\end{table}

Dies liegt vermutlich daran, dass beide Modell auf weltweiten Meteoparametern basieren. Wie Infinity jedoch die «berechnete»
Version erstellt ist weiterhin unklar. Gemäss Empfehlung wird für die finalen GNSS-Koordinaten das Modell «Berechnet» verwendet.
(\todo{Effektiv Koord neu rechnen oder Satz anpassen})


\subsubsection{Einfluss der externen Referenzstation}
Die vier lokalen Referenzstationen werden zum Verlgeich sowohl zur Agnes-Station HABG, als auch zu einer 
virtuellen Referenzstation (VRS) berechnet. Die VRS wird hierbei auf dem Punkt EX9500 auf einer Höhe von 650 m.ü.M. gewählt. 
Die Basislinien zur VRS weisen für alle Stationen deutlich mehr und grössere Ausreisser im GDOP über die gesamte Messzeit auf. 
Während der GDOP der Basislinien zur HABG 4.5 nicht überschreitet, steigt dieser stellenweise auf über 15 für Basislinien zur VRS.
Die Koordinatendifferenzen zwischen den Berechnungen zur HABG und zur VRS betragen im Mittel 
\todo{Koodrdinatendifferenzen und RMSE berechnen HABG/VRS} mit einem RMSE von \todo{einfügen}. 
Da abgesehen der HABG alle umliegenden Swipos-Referenzstationen 40 km und weiter entfernt und meist auf Höhen 
ausserhalb der vertikalen Projektausdehnung liegen, ist durch das Verwenden der VRS keine Genauigkeitssteigerung zu 
erwarten.

Welche Referenzstation in den Vorjahren verwendet wurde, ist nicht dokumentiert. Es ist jedoch naheliegend, dass eine 
AGNES-Station verwendet wurde, da diese näher liegt als die umliegenden Refnet-Stationen. Zudem gab es zur Zeit 
der ersten GNSS-Messungen am Schwanderbärgli vermutlich ausschliesslich AGNES-Stationen in der Schweiz.

\todo{evtl. Abbildung GDOP Werte (oder in Anhang)}


\subsection{Resultate}
\label{sec:resultate}

\subsection{Diskussion der Resultate}
\label{sec:diskussion der resultate}

